% ==============================================================================
% MÓDULO: CAPÍTULO I - EL PROBLEMA DE INVESTIGACIÓN
% Archivo: subcapitulos/sec_01_problema.tex
% Enfoque: 100% Calidad de Atención (Donabedian y Watson)
% Cátedra: Proyecto de Investigación en Salud (EN 486) - UNSCH
% ==============================================================================
\section{EL PROBLEMA DE INVESTIGACIÓN}
\label{sec:cap1_problema}

\subsection{Contextualización del problema de investigación}
\label{sec:1.1_contextualizacion_problema}

En el marco epistemológico de la investigación cualitativa y la guía oficial para formular proyectos con enfoque cualitativo de la Escuela Profesional de Enfermería (UNSCH), el planteamiento del problema adopta una \textbf{naturaleza puramente inductiva y fenomenológica}: no parte del aislamiento artificial de variables numéricas ni del contraste de hipótesis cerradas, sino del sentir, pensar y reflexionar sobre la trayectoria vivencial en el escenario natural donde acontece el encuentro del cuidado humano \cite{sampieri2018}. Como señalan Husserl \cite{husserl1949} y van Manen \cite{vanmanen1990}, contextualizar un problema fenomenológico exige capturar la esencia, las emociones, los sentires corpóreos y los significados intersubjetivos que las personas otorgan a la calidad del cuidado recibido en el momento del acto asistencial.

Para tal fin, la metodología demanda un \textbf{contacto directo y vivencial con la realidad estudiada en su escenario cotidiano}. En el presente estudio, el origen de la indagación se sustenta en la observación reflexiva y la experiencia acumulada por la investigadora durante sus prácticas preprofesionales de la Escuela Profesional de Enfermería de la Universidad Nacional de San Cristóbal de Huamanga (UNSCH) en los consultorios del Centro de Salud Belén, dialogando con los usuarios dolientes y presenciando sus incertidumbres y sentires.

En estricta observancia del canon metodológico del Dr. Manglio Aguirre Andrade y el protocolo cualitativo de referencia, la contextualización del problema de calidad de atención se articula a través del \textbf{recorrido inverso}, partiendo del ámbito específico hacia el general:

\subsubsection*{1. Ámbito Local (Centro de Salud Belén, Ayacucho)}
En el plano local, el estudio se centra de manera estricta y exclusiva en el Centro de Salud Belén (Categoría I-3, Microred Huamanga), establecimiento cabecera que atiende a una población adscrita de 18,450 habitantes en el distrito de Ayacucho, provincia de Huamanga, a 2,760 metros sobre el nivel del mar. La investigación se focaliza exclusivamente en la dinámica asistencial de Belén, en vinculación directa con la \textbf{Prioridad Regional 10 de Investigación en Salud de Ayacucho 2025–2030}, aprobada por la Dirección Regional de Salud (DIRESA) de Ayacucho \cite{diresa2025}: \textit{“Organización, gestión, calidad y accesibilidad de los servicios de salud”}.

En este escenario aflora una profunda distancia entre las normas programáticas institucionales y la vivencia sentida del usuario andino frente a la calidad del cuidado:
\begin{itemize}
    \item \textbf{Vulnerabilidad socioeconómica y necesidad de contención afectiva:} El 68.2\% de la población usuaria se encuentra en situación de pobreza según el SISFOH y más del 94.5\% depende del Seguro Integral de Salud (SIS). Para estas familias, acudir a los consultorios de Belén representa un momento de vulnerabilidad extrema donde demandan no solo una prescripción médica o un procedimiento de enfermería, sino contención afectiva y comprensión empática frente a su precaria situación de vida.
    \item \textbf{Identidad lingüística y códigos andinos de calidad:} El 72.0\% de los usuarios presenta bilingüismo activo quechua chanka-castellano. Para esta población, la calidad de la atención no se concibe como un trámite administrativo frío, sino como un acto de reciprocidad y respeto que exige la acogida hospitalaria (\textit{allin chaskiy}), el saludo cálido (\textit{napaykuy}) y el respeto mutuo (\textit{respetanakuy}), conforme a los hallazgos de Aguirre-Andrade y colaboradores \cite{aguirre2020}.
    \item \textbf{La vivencia de la desatención y la prisa protocolar:} Cuando el personal sanitario —condicionado por la sobrecarga laboral o la digitación continua en el computador— no establece contacto visual, interrumpe el relato del paciente o muestra indiferencia ante el dolor físico (\textit{nanay}), el usuario experimenta una profunda insatisfacción afectiva y desamparo, replegándose en el silencio defensivo (\textit{upallay}). Asimismo, las consultas abreviadas de escasos minutos dejan en las madres de familia en CRED y en las personas adultas mayores dudas e incertidumbres no resueltas.
\end{itemize}

\subsubsection*{2. Ámbito Nacional (Perú)}
En el ámbito nacional, la realidad observada en Belén refleja los nudos críticos del primer nivel de atención en el Perú (categorías I-1 a I-4), donde se concentra más del 78\% de las atenciones asistenciales de todo el país. La Política Nacional de Calidad en Salud del Ministerio de Salud \cite{minsacalidad2021} reconoce que la persistencia de barreras de trato interpersonal, frialdad burocrática y falta de pertinencia cultural debilita gravemente la capacidad resolutiva del sistema público. Esta situación motivó al Estado peruano a priorizar la investigación en servicios sanitarios a través de la \textbf{Línea 10 de Investigación en Salud al 2030}: \textit{“Sistemas y servicios de salud, acceso y cobertura universal”} (Resolución Ministerial N° 424-2025/MINSA) \cite{minsainv2025}, buscando dignificar la atención sanitaria en todos los establecimientos de atención primaria de la nación.

\subsubsection*{3. Ámbito Mundial (Internacional o Global)}
A escala mundial, los sistemas de salud enfrentan una severa crisis de deshumanización asistencial caracterizada por la tecnificación desmedida del acto clínico y la sobrecarga burocrática en el primer nivel de atención, según lo advertido por la Organización Mundial de la Salud (OMS), la OCDE y el Banco Mundial \cite{who2018}. El predominio de un paradigma biomédico mecanicista ha reducido la evaluación de los servicios a meros indicadores cuantitativos de producción, invisibilizando la vivencia afectiva del paciente, el sufrimiento somático y los determinantes socioculturales de la salud. De acuerdo con Avedis Donabedian \cite{donabedian1980, donabedian2005} y Jean Watson \cite{watson2008}, la calidad del cuidado no reside únicamente en la pericia instrumental de un procedimiento, sino medularmente en el proceso interpersonal: la calidez de la mirada, la capacidad de escucha compasiva, el respeto a la dignidad y la reciprocidad humanizada entre el profesional de salud y el ser cuidado.

El problema medular y objeto de estudio de la investigación radica, por consiguiente, en develar y comprender cómo los usuarios significan, valoran y experimentan la calidad de atención y el cuidado de enfermería en los consultorios del Centro de Salud Belén desde la profundidad de su mundo de la vida (\textit{Lebenswelt}).

---

\subsection{Formulación del Problema (Preguntas Norteadoras)}
\label{sec:1.2_formulacion_problema}

En coherencia con los preceptos del diseño fenomenológico de Hernández-Sampieri y Mendoza \cite{sampieri2018} y el marco teórico de la calidad interpersonal de Donabedian \cite{donabedian2005} y Watson \cite{watson2008}, se formulan las preguntas norteadoras de la investigación:

\subsubsection{Problema General (Pregunta Principal)}
\label{sec:1.2.1_preg_gral}
¿Cuál es el significado, estructura y esencia de la experiencia vivida por los usuarios respecto a la calidad de atención recibida en los servicios de salud del Centro de Salud Belén, Ayacucho 2026?

\subsubsection{Problemas Específicos (Preguntas Específicas)}
\label{sec:1.2.2_preg_esp}
\begin{enumerate}
    \item ¿Cómo experimentan los usuarios la \textbf{dimensión relacional (relación humana vivida / \textit{Mitwelt})} respecto a la calidez, empatía, escucha activa y comunicación intercultural en quechua chanka por parte del personal de salud?
    \item ¿De qué manera vivencian los usuarios la \textbf{dimensión temporal (tiempo vivido / \textit{Lebenszeit})} respecto a la puntualidad, oportunidad y tiempo efectivo dedicado por el profesional durante la consulta clínica?
    \item ¿Cómo vivencian los usuarios la \textbf{dimensión espacial (espacio vivido / \textit{Lebensraum})} respecto a la privacidad, confidencialidad, comodidad física y dignidad ambiental de los consultorios?
    \item ¿Cuáles son las vivencias de los usuarios vinculadas a la \textbf{dimensión corporal (cuerpo vivido / \textit{Leib})}, tales como el respeto al pudor físico, la contención del sufrimiento y el alivio del dolor (\textit{nanay})?
    \item ¿Cuál es la \textbf{esencia compartida y las visiones divergentes} que configuran el significado atribuido por los usuarios a una atención de salud con calidad y calidez en el Centro de Salud Belén?
\end{enumerate}

---

\subsection{Objetivos de la Investigación}
\label{sec:1.3_objetivos}

Guardando una estricta simetría biunívoca (1:1) con las preguntas norteadoras, se establecen los siguientes objetivos de investigación:

\subsubsection{Objetivo General}
\label{sec:1.3.1_obj_gral}
Comprender el significado, estructura y esencia de la experiencia vivida por los usuarios respecto a la calidad de atención recibida en el Centro de Salud Belén, Ayacucho 2026.

\subsubsection{Objetivos Específicos}
\label{sec:1.3.2_obj_esp}
\begin{enumerate}
    \item Comprender y analizar la experiencia de los usuarios en la \textbf{dimensión relacional (relación humana vivida / \textit{Mitwelt})} respecto a la calidez interpersonal, el trato digno y la comunicación en lengua quechua chanka brindada por el equipo de salud.
    \item Explorar y comprender las vivencias de los usuarios en la \textbf{dimensión temporal (tiempo vivido / \textit{Lebenszeit})} en relación con la puntualidad, prontitud y duración efectiva del acto del cuidado asistencial en la consulta.
    \item Describir y comprender la experiencia de los usuarios en la \textbf{dimensión espacial (espacio vivido / \textit{Lebensraum})} en cuanto al confort, reserva de la intimidad y privacidad física durante el cuidado en consultorios.
    \item Identificar y comprender las vivencias de la \textbf{dimensión corporal (cuerpo vivido / \textit{Leib})}, analizando el respeto al pudor físico del paciente y la respuesta humanizada ante el dolor y sufrimiento somático (\textit{nanay}).
    \item Sintetizar la \textbf{estructura esencial y las divergencias fenomenológicas} que componen la calidad de atención de salud en el Centro de Salud Belén.
\end{enumerate}

---

\subsection{Viabilidad de la Investigación y Consideraciones Bioéticas}
\label{sec:1.4_viabilidad_bioetica}

El estudio es plenamente viable en términos temporales, técnicos y socioculturales, al contar con la apertura de los usuarios y el respaldo institucional de la cátedra de EN 486. La investigación se sujeta rigurosamente a los principios de la Declaración de Helsinki \cite{helsinki2013} y las pautas éticas internacionales del CIOMS \cite{cioms2016}:
\begin{itemize}
    \item \textbf{Autonomía y Consentimiento Informado:} Se aplicará un consentimiento informado bilingüe detallado, garantizando que la participación sea voluntaria y revocable en cualquier momento.
    \item \textbf{Confidencialidad y Anonimato:} Las grabaciones y transcripciones serán anonimizadas mediante códigos alfanuméricos confidenciales (INF-01 a INF-18).
    \item \textbf{No maleficencia y beneficencia:} La indagación no genera riesgos físicos ni psicológicos, y sus aportes servirán para promover mejoras directas en la humanización del cuidado en el Centro de Salud Belén.
\end{itemize}

\newpage
